\section{System Design}
\label{sec:design}

Konstruct reads a repository and the warehouses its dbt profiles name,
and writes \cmd{konstruct-out/graph.json}, \cmd{REPORT.md}, \cmd{graph.html} and
\cmd{ontology.yml}. Three rules shape it. It never imports or executes user orchestration code:
DAG, Dagster and PySpark files are parsed with \cmd{ast}, dbt Jinja is rendered in a sandbox, and
\cmd{konstruct verify} runs a model's SQL only on copies of its inputs or inside a transaction that
is rolled back; on DuckDB copies it turns DuckDB's external access and extension autoloading off, so a
model that reads files or URLs directly is skipped. The one exception is dbt itself: by default, when a dbt project's artefacts are
missing or stale, \cmd{konstruct build} runs \cmd{dbt docs generate}, which runs the project's
macros and on-run hooks with its own profile, and those can write to the warehouse; with
\cmd{--no-dbt}, the GitHub Action's default, dbt never runs. A repository's own warehouse
connections, the paths it names outside the checkout and running dbt are used only once the user
has reviewed the project and run \cmd{konstruct trust}, which records a digest of those settings
outside the repository, so that any change to them needs trust again. Apart from that, it reads warehouses
for metadata only, makes no network requests of its own beyond the warehouses it is configured to
read, and collects no telemetry; when it runs dbt, it turns dbt's anonymous usage statistics off.
And everything it claims about lineage is tested, including against data (\cref{sec:verify}).

\subsection{The Graph Model}
\label{sec:graph}

The graph is a typed, directed multigraph $G = (V, E)$. Each node $v \in V$ has a kind
$\mathit{kind}(v)$ among \emph{source}, \emph{table}, \emph{model}, \emph{seed}, \emph{snapshot},
\emph{column}, \emph{dag}, \emph{task}, \emph{test}, \emph{exposure} and \emph{term}; each edge
$(u, v, k) \in E$ has a kind $k$ among \emph{contains}, \emph{feeds}, \emph{derived\_from},
\emph{runs}, \emph{precedes}, \emph{tests}, \emph{defines} and \emph{relates\_to}, and carries
attributes. Let $\R \subseteq V$ be the relations (sources, tables, models, seeds and snapshots)
and $C \subseteq V$ the columns, each with a parent relation $p(c) \in \R$. Two relations are at
the heart of the model:
\begin{align}
  \D &= \{(c, c', \tau) : c \text{ is derived from } c'\} \subseteq C \times C \times \Theta, \label{eq:lineage}\\
  \mathcal{F} &= \{(c', r) : r \text{ joins, filters or groups on } c'\} \subseteq C \times \R. \label{eq:refs}
\end{align}
The transform $\tau \in \Theta = \{\textit{direct}, \textit{transformed}\}$ is \emph{direct} when the path from $c'$ to $c$ is made only of column
references and renames, and \emph{transformed} otherwise. $\D$ says where a value comes from;
$\mathcal{F}$ says which columns decide the rows a relation returns, and is stored as
\emph{feeds} edges with the role \emph{reference}. Keeping the two apart matters: a filter column
changes which rows a model returns, but marking every downstream column as derived from it would
bury the answer to ``what does this column depend on''.

Each relation has a canonical key $\rel(r)$, its lower-case \cmd{catalog.schema.name}. Nodes from
different tools that describe one physical relation (a dbt model, the warehouse table it
materialises, the target of an Airflow SQL task, an Iceberg table) share a key, and SQL that names
the key resolves to one of them by a fixed priority, so a chain of derivations crosses tools
without any configuration. \cmd{impact} and \cmd{trace} are the forward and backward closures
over $\D$, $\mathcal{F}$ and the task edges, reported with the transform at each step and the
business entities the walk reaches; \cmd{path} is a shortest path between two nodes in either
direction.

\subsection{SQL Column Lineage and Canonical Expressions}
\label{sec:sql}

For each statement, Konstruct qualifies every column against the columns it knows: those from
the dbt catalog or the warehouse, and those of upstream models analysed earlier in dependency
order, so \cmd{SELECT *} expands. It then follows each output column through CTEs, subqueries,
aliases, unions, window functions and nested fields to physical columns with sqlglot's
lineage~\citep{sqlglot}, in the dialect of the dbt adapter, one of the \sqlDialects{} dialects that adapters map to.
\cmd{CREATE ... AS}, \cmd{INSERT ... SELECT}, \cmd{MERGE}, \cmd{UPDATE}, \cmd{DELETE} and
\cmd{COPY} give the relation they write.

Alongside the edges, Konstruct keeps each output column's \emph{canonical expression}: its SQL
with every column replaced, recursively, by the canonical expression of what it reads, until only
physical columns remain. For a lineage node $n$,
\begin{equation}
  \kappa(n) =
  \begin{cases}
    \rel(t).a & \text{column $a$ of table $t$},\\
    \mathrm{op}\big(\mathrm{sort}\,\kappa(\mathrm{ch}(n))\big) & \text{set operation},\\
    e_n\big[x \mapsto \langle \kappa(\mathrm{ch}_x(n)) \rangle\big] & \text{otherwise},
  \end{cases}
  \label{eq:canonical}
\end{equation}
where $e_n$ is the node's expression, $\mathrm{ch}_x(n)$ the child that column $x$ resolves to,
and $\langle\cdot\rangle$ brackets an inlined expression unless it is a name. The join, filter,
grouping, \cmd{HAVING} and \cmd{QUALIFY} conditions of each scope are kept in the same form.
Because CTE names, aliases, quoting, formatting and column order disappear in $\kappa$, two
versions of a model that differ only in those have equal canonical expressions. \cmd{konstruct
diff} relies on this: it reports a change only when a column's canonical expression or a
condition changes, and a removed column as breaking only when something downstream still reads
it. A model whose Jinja has not been compiled is rendered by Konstruct's own sandboxed renderer,
with \cmd{ref}, \cmd{source}, \cmd{var}, the project's and packages' macros and
\cmd{adapter.dispatch}; macros that would query the warehouse return empty results.

\subsection{Formula Composition}
\label{sec:compose}

A canonical expression stops at the boundary of a model. \cmd{konstruct explain} composes it
through every model and task upstream. For a column $c$ with canonical expression $e(c)$ and
inputs $\inputs(c) = \{c' : (c, c', \tau) \in \D\}$, the composed formula is
\begin{equation}
  \varphi(c) =
  \begin{cases}
    c & \text{if } \inputs(c) = \emptyset,\\
    e(c)\big[c' \mapsto \langle \varphi(c') \rangle\big]_{c' \in \inputs(c)} & \text{otherwise},
  \end{cases}
  \label{eq:compose}
\end{equation}
a source column standing for itself. In the bundled example, \cmd{customer\_ltv.lifetime\_value} reads as
\cmd{\exLtvFormula{}}, across an Airflow SQL task, a dbt mart and a staging model. Composition is
memoised, so each column is composed once however many columns read it; a cycle, a depth beyond
\maxDepth{} columns or a formula longer than \maxComposed{} characters stops substitution and
leaves the column as a reference, and the formula is then marked incomplete. The budget is checked
before substituting, since a formula can grow quickly when a column reads several inputs that
each compose to long formulas.

$\varphi$ describes values, not rows. Substitution is exact for the expression, but the joins,
filters and grouping of every step decide which rows reach it, and folding them into a single
expression would be wrong for aggregates. \cmd{explain} therefore lists each step with its
conditions beside the formula, and a test evaluates composed formulas of chains without joins,
filters or aggregation directly on the source tables and requires them to give the final model's
values row for row.

\paragraph{Composition as text with operand-sensitive contexts}
The viewer shows every column's composed formula, so a build composes all of them. Building each
$\varphi(c)$ as a tree copies and prints ever larger trees. Konstruct instead prints each column's
own expression once, with each input replaced by a placeholder token, and splices the composed
texts of the inputs into the printed string, adding the brackets the composed tree would keep.
This is only correct where the printer's output for a node does not depend on what its operand is:
BigQuery prints \cmd{MOD} differently by operand, T-SQL rewrites a non-boolean condition as a
comparison with zero, and DuckDB's string functions and \cmd{INTERVAL} and \cmd{ANY} look at their
arguments. Konstruct keeps an explicit set $T$ of \textContexts{} node kinds whose printing is
operand-insensitive in every dialect (arithmetic, comparisons, \cmd{COALESCE}, \cmd{CASE}, the
core aggregates, \cmd{CAST}, \cmd{IN}, \cmd{LIKE}, window specifications, field access and a few
others), and uses text only when every input sits in such a context:

\begin{algorithm}[t]
\caption{Composing a column's formula as text}
\label{alg:text}
\begin{algorithmic}[1]
\Require column $c$, memo $M$
\If{$c \in M$} \Return $M[c]$ \EndIf
\State $t \gets$ parse$(e(c))$;\ $S \gets \emptyset$
\For{each column reference $x$ in $t$ resolving to $c' \in \inputs(c)$}
  \State $s \gets$ \Call{ComposeText}{$c'$}
  \If{the parent of $x$, past brackets, is not in $T$} \Return \Call{ComposeTree}{$c$} \EndIf
  \State replace $x$ by a fresh token $k$;\ $S[k] \gets s$
\EndFor
\State decide, for each $k$, whether $S[k]$ needs brackets in its parent; drop redundant ones
\State $p \gets$ print$(t)$ in the dialect
\If{a token is printed other than exactly once} \Return \Call{ComposeTree}{$c$} \EndIf
\State $M[c] \gets$ $p$ with each token $k$ replaced by $S[k]$
\State \Return $M[c]$
\end{algorithmic}
\end{algorithm}

\noindent The invariant is that \cref{alg:text} returns exactly the string that printing the tidied
tree of \cref{eq:compose} would. It holds by construction where it uses text: a node in $T$ prints
its operands' strings unchanged, and the bracket decision uses the kind of node an input's
formula starts with (before and after its own brackets are tidied), which is what the tree printer
looks at. Elsewhere it falls back to the tree. The test suite checks the invariant on every column
of a built project and on each input shape in each context in several dialects.

\subsection{One Chain Across Tools}
\label{sec:chains}

Every input contributes nodes and edges to the same graph, keyed so that the same relation is one
node:
\begin{itemize}
  \item \textbf{dbt}: models, seeds, snapshots, sources, tests, constraints and exposures from the
  manifest and catalog, with compiled SQL for lineage, or SQL rendered without dbt.
  \item \textbf{SQL}: standalone files and warehouse view definitions.
  \item \textbf{Airflow}: DAGs, TaskFlow, TaskGroups and every dependency syntax, parsed with
  \cmd{ast}; dbt commands in Bash, Cosmos and dbt operators resolved to models with dbt's selection
  syntax; inline and file SQL with column lineage; inlet and outlet assets.
  \item \textbf{Dagster}: assets, multi-assets and dbt assets with their selections, op jobs
  (nested graphs, aliases, fan-in, dynamic mapping), schedules and sensors, with column lineage
  from the SQL an asset or op runs.
  \item \textbf{PySpark}: the SQL a job passes to a Spark session, in the Spark dialect.
  \item \textbf{Dashboards}: dbt exposures, Looker views field by field from LookML, and Power BI
  semantic models from TMDL, with DAX references resolved across tables.
  \item \textbf{Loaders}: the tables Fivetran and Airbyte land, as sync tasks per schema, and
  column lineage from each Fivetran source system when its Platform Connector's tables are in the
  warehouse.
  \item \textbf{Warehouses}: DuckDB, Postgres, Snowflake, BigQuery, Databricks, Trino and Iceberg
  through any PyIceberg catalogue, read for tables, columns, types, comments, keys and views.
  \item \textbf{Meaning}: \cmd{definitions.yml}, SKOS and OWL vocabularies~\citep{skos,owl2},
  Apache Ossie models, and column usage from the query history of \usageWarehouses{} warehouses.
\end{itemize}
A column computed in dbt, copied into a reporting table by an Airflow SQL task, written by a
Dagster asset and read by a Looker field is therefore one chain, and so is a chain across several
repositories that share a warehouse.

\subsection{Builds, Caching and Parallelism}
\label{sec:build}

Analysing SQL dominates a build. Each analysis is cached under a content key; for a dbt model $m$,
\begin{equation}
  k(m) = \mathrm{SHA256}\big(\textit{sql}_m,\ d,\ \rel(m),\ S(\mathrm{up}(m)),\ \textit{checksum}_m\big),
  \label{eq:cache}
\end{equation}
where $d$ is the dialect and $S(\mathrm{up}(m))$ the columns of the models $m$ reads. Because the
key holds the upstream schemas, an edit that changes a model's columns invalidates exactly the
analyses that could see the change. \cmd{konstruct build --update} reuses every other analysis
and recomputes only the formulas downstream of the change, but it still reads every input, checks
every key and writes every output, so its cost grows with the project as well as with the change
(\cref{sec:eval-performance}). A cold build analyses dbt models a level of the DAG at a time in
worker processes, one per CPU up to \maxWorkers{}, when at least \parallelMinimum{} models are
ready, with results identical to analysing them in order.

\subsection{The Ontology}
\label{sec:ontology}

Konstruct derives an ontology from the graph, written to \cmd{ontology.yml} on every build and
shown by \cmd{konstruct ontology}. Most of it is read from code and tests rather than declared:

\paragraph{Entities} A column $k$ is an \emph{entity key} when its name is identifier-like
(\cmd{customer\_id}, \cmd{order\_key}, \cmd{id}) and a dbt \cmd{unique} test vouches for it, or
when it is a single-column primary key (a dbt constraint, a warehouse primary key, an Iceberg
identifier field, an Ossie primary key). The entity is named after the key, or after the relation
for a bare \cmd{id}, and every relation keyed by it holds one row per instance. A term of kind
\emph{entity} names one explicitly, and a measure's grain (``one row per store'') finds the key of
the entity it names.

\paragraph{Relationships} A dbt \cmd{relationships} test, an Ossie relationship, or a join that
equates a column of one entity's relation with another entity's key gives a relationship. Its
cardinality is claimed only on evidence:
\begin{equation}
  \mathrm{card}(a \to b) =
  \begin{cases}
    \text{one-to-one} & u(a) \wedge u(b),\\
    \text{many-to-one} & u(b) \wedge \neg u(a),\\
    \text{unknown} & \text{otherwise},
  \end{cases}
\end{equation}
where $u(x)$ holds when a test or key shows $x$ unique.

\paragraph{Measures} A business term of kind \emph{measure}, an Ossie metric, or a term bound to
computed columns is a measure. Each bound column carries $\varphi(c)$, the sources it reaches,
how many lineage edges on the way \cmd{verify} confirmed, and its usage when query history is read.

\paragraph{Rules} The build holds the ontology to \ontologyRules{} rules; a violation is a warning,
so \cmd{konstruct build --strict} fails on it. The strongest is \emph{one formula per measure}:
every column bound to a measure must compute it the same way, compared as \cmd{konstruct
conflicts} compares columns, by composed formula, sources, filters and joins on the route back to
source, and the aggregation at the top. The others require entity keys to be unique, entity names
to be distinct, bound columns to exist and approved terms to have an owner and an approval date.

\subsection{Apache Ossie Round Trips}
\label{sec:ossie}

Konstruct reads every \cmd{*.ossie.yaml} file, in the released shape, with models in a list, and the newer draft, with one model at the root:
a dataset becomes the relation its source names, a field the source columns its expression reads,
a metric a measure, a primary key an entity key and a relationship a relationship. Anything that
does not resolve is a build warning, including a model-scoped metric that names a field without
its dataset, which the specification does not allow. \cmd{konstruct export ossie} writes the
ontology as one model. Writing is \emph{authored-preserving}: a dataset read from an Ossie model is
written back as authored, and so is each metric whose datasets keep their names, with a
\cmd{KONSTRUCT} custom extension added that holds the metric's status, owner, approval date,
entity, and each bound column's formula, sources, verified edges and usage. Every other metric is
derived from its composed formula, with the portable expression in ANSI SQL and the warehouse
dialect's expression beside it. If $R$ reads a model and $W$ writes one, the property tested on
Ossie's own portability ladder is
\begin{equation}
  W(R(M)) = M \oplus x
  \label{eq:roundtrip}
\end{equation}
for every model $M$ whose metrics name their fields as the specification requires, where $\oplus\, x$ adds only the
extension. \Cref{sec:eval-ossie} grades each element.

\subsection{Ontology Packs}
\label{sec:packs}

Many teams load the same sources. An ontology pack is a definitions file for one source schema
with a match block; Konstruct ships \packCount{} (\cref{tab:packs}). A schema $\sigma$ matches a
pack $P$ when it has every table $P$ requires with its required columns, and there is evidence
that it comes from the source: its name fits one of the pack's patterns, or a loader landed its
tables. The second condition keeps an application's own \cmd{orders} and \cmd{customers} tables
from being taken for Shopify's. A match binds the pack's terms to the schema's columns as drafts,
leaving out terms whose tables or columns are missing, and the project's own terms win over a
pack's.

\begin{table}[t]
\centering
\caption{The ontology packs Konstruct ships: the tables a schema must have, and the entity,
measure and other terms each binds.}
\label{tab:packs}
\small
\begin{tabular}{@{}lrrrr@{}}
\toprule
Pack & Tables & Entities & Measures & Terms\\
\midrule
\packsTable
\bottomrule
\end{tabular}
\end{table}

\subsection{Interfaces and the Agent Gate}
\label{sec:agents}

The same questions are answered everywhere. The command line has \cliCommands{} commands:
\cmd{konstruct init}, \cmd{konstruct build}, \cmd{konstruct status} and \cmd{konstruct open} set up,
build, check and show the graph, and \cmd{konstruct trust} trusts a project's own settings;
\cmd{konstruct impact}, \cmd{konstruct trace}, \cmd{konstruct
path}, \cmd{konstruct search}, \cmd{konstruct describe}, \cmd{konstruct list} and \cmd{konstruct
query} walk and find it; \cmd{konstruct explain} and \cmd{konstruct history} give a column's
formula and when it changed; \cmd{konstruct verify} checks lineage against data, and with
\cmd{--mutate} whether a model's dbt tests fail on small changes to its SQL, and with
\cmd{--why} the source rows behind one value with each row's share;
\cmd{konstruct diff} and \cmd{konstruct gate} review changes; \cmd{konstruct receipts} records
what an answer rests on (each column's composed formula, row conditions, sources, bound terms and
verification state) and later reports the answers whose basis has changed;
\cmd{konstruct define},
\cmd{konstruct ontology} and \cmd{konstruct conflicts} answer what things mean, and
\cmd{konstruct check-sql} checks a query against them without running it;
\cmd{konstruct export} writes OpenLineage events (\cmd{openlineage}), Mermaid (\cmd{mermaid}),
Graphviz (\cmd{dot}), CSV (\cmd{csv}) and Apache Ossie (\cmd{ossie}); \cmd{konstruct plugins}
lists the installed plugins, which add inputs, export formats, gate checks, ontology packs and
MCP tools, and whether each runs; and \cmd{konstruct install}, \cmd{konstruct mcp} and
\cmd{konstruct hook} connect coding agents. The Python API, \cmd{konstruct.Project}, returns
typed results for each question, and \cmd{graph.html} is one offline file with the lineage,
formulas, terms, verification results and usage. \cmd{konstruct serve} sends the same viewer from
a local server, which reloads open pages when a build replaces the graph, beside a read-only HTTP
API, described by an OpenAPI document, whose answers come from the Python API's own functions.
With \cmd{--meaning}, \cmd{konstruct search} also finds a node by what it is rather than what it
is called: identifiers are read as words, with the abbreviations and synonyms of data work, and
each node is scored, with no trained model, on the words, formula tokens and lineage it shares
with the question.

\paragraph{Grounding agents} \cmd{konstruct mcp} runs a Model Context Protocol server with
\mcpTools{} tools: \cmd{impact}, \cmd{trace}, \cmd{path}, \cmd{search}, \cmd{describe},
\cmd{query}, \cmd{explain}, \cmd{history}, \cmd{define}, \cmd{conflicts}, \cmd{ontology},
\cmd{diff}, \cmd{gate}, \cmd{check\_sql}, \cmd{check\_history}, \cmd{mutation\_test} and
\cmd{why}. Answers are the same structured results the Python API returns, so an
agent learns a column's formula, its sources, which of its edges the data confirmed, the approved
definition of the business term it holds, its owner and the entities it touches, before it writes
a query that joins or aggregates. The server rereads the graph whenever it changes, and
\cmd{konstruct install} writes skills for Claude Code, Cursor, Codex, GitHub Copilot, Gemini CLI and
Windsurf that teach the agent whole workflows, such as a safe change, a review or a correct query,
with the commands each step needs.

\paragraph{Checking a query} An agent that answers a business question writes SQL, and SQL that
runs can still be wrong. \cmd{konstruct check-sql} reads a query without running it and composes
each aggregated output column down to source columns, with the filters on its way, as the graph's
columns are composed. It reports a measure that leaves out a filter its approved definition
requires or computes it differently, a column of a retired term, a sum, count or average over a
join that repeats the rows it aggregates, judged from the keys, tests and relationships of the
ontology, and names the project does not have. Each finding carries its fix, and where a fix
changes only the measure the result carries the corrected query. A match by name alone is never
an error, since an agent told that a correct query is wrong will break it. With
\cmd{--history}, the same checks run over a warehouse's query history and count the queries,
runs and people that compute an approved measure differently from its definition.

\paragraph{The gate} An agent that changes a data stack should prove the change is safe before it
proposes it. \cmd{konstruct gate} compares the working tree with a base, as \cmd{diff} does, and
fails when the change (a) removes a column that a model, task or dashboard still reads; (b)
changes what an approved term means, by changing its wording, columns, filters or grain, or the
canonical expression or row conditions of a column it names or of any column upstream of one; or
(c) adds a column that computes the same business idea as an existing one in a different way, as
\cmd{conflicts} finds, unless the two are bound to different terms. Each finding says what to do,
so an agent can fix the change and run the gate again. \cmd{konstruct hook} runs the gate as a
Claude Code hook after an agent edits SQL or dbt YAML, from files alone, so the findings arrive
next to the edit. The GitHub Action runs the same check on a
pull request when it is configured with \cmd{gate: true}.
